\pdfbookmark[0]{Résumé}{Résumé}
\chapter*{Résumé}
\label{chap:abstract-fr}

Ce projet de fin d'année, réalisé chez Huawei Maroc, présente la conception et la mise en œuvre d'une plateforme intelligente pour l'analyse de la migration de GPU vers NPU et l'aide à la décision en matière d'infrastructure d'IA.

Le projet répond au défi du déploiement de modèles d'IA dans des environnements matériels hétérogènes, particulièrement la transition des écosystèmes GPU basés sur CUDA vers les NPU Huawei Ascend. Le système proposé offre un cadre unifié pour analyser les charges de travail des modèles, évaluer la compatibilité matérielle, estimer la faisabilité de la migration et recommander des stratégies de déploiement optimales.

La plateforme repose sur une architecture modulaire composée de services d'analyse de modèles, de vérification de compatibilité, de benchmarking et de génération de recommandations. Elle introduit une couche d'abstraction matérielle permettant l'évaluation sur les infrastructures CPU, GPU et les futures infrastructures NPU.

La mise en œuvre suit les pratiques modernes de génie logiciel, incluant le déploiement conteneurisé et l'orchestration de services via API. Le benchmarking est utilisé comme mécanisme de validation pour soutenir les décisions de migration plutôt que comme un objectif primaire.

Le système démontre une approche structurée de la prise de décision en matière d'infrastructure d'IA en reliant les écosystèmes CUDA et Huawei Ascend, fournissant une base pour de futures optimisations et des stratégies de migration automatisées.

\vspace{0.5cm}
\noindent\textbf{Mots-clés :} migration GPU-NPU, Huawei Ascend, CUDA, infrastructure d'IA, abstraction matérielle, compatibilité de modèle, évaluation de performance
